\section{Height functions}

Let $S=(K,(\Omega,\mathcal A,\nu),\phi)$ be an adelic structure and
let $I$ be a finite nonempty index set. We first construct heights over
$K$ itself; applying the construction to $S\otimes_K K^{\mathrm{ac}}$
gives Definition 3.5.1 of \cite{chen-moriwaki}. Properness is required
for descent to projective space, but not for integrability.

\begin{definition}[Local maximum and its logarithm]\label{def:local-height}
\uses{def:adelic-structure}
\lean{ArakelovGeometry.AdelicStructure.localHeight,ArakelovGeometry.AdelicStructure.localHeight_zero,ArakelovGeometry.AdelicStructure.eval_le_localHeight,ArakelovGeometry.AdelicStructure.localHeight_pos}
\leanok
For $a=(a_i)_{i\in I}\in K^I$ set
$M_S(a,\omega)=\max_{i\in I}|a_i|_\omega$ and
$\ell_S(a,\omega)=\log M_S(a,\omega)$ for $a\ne0$.
For the zero tuple put $\ell_S(0,\omega)=0$.
\end{definition}

\begin{lemma}[Measurability and integrability of local heights]
\label{lem:local-height-integrable}
\uses{def:local-height}
\lean{ArakelovGeometry.AdelicStructure.measurable_log_eval,ArakelovGeometry.AdelicStructure.integrable_log_eval,ArakelovGeometry.AdelicStructure.measurable_eval,ArakelovGeometry.AdelicStructure.measurable_localHeight,ArakelovGeometry.AdelicStructure.integrable_log_localHeight}
\leanok
For every $a\in K^I$, the function $\ell_S(a,\cdot)$ is measurable
and integrable.
\end{lemma}
\begin{proof}
For $b\ne0$, $|b|_\omega=\exp(\log|b|_\omega)$ is measurable;
for $b=0$ it is constant. Thus the finite maximum $M_S(a,\cdot)$
is measurable. For $a\ne0$, any nonzero coordinate shows that this
maximum is positive, so its logarithm is measurable. At each $\omega$
choose an index attaining the maximum. Its coordinate is nonzero and
therefore
\[
 |\ell_S(a,\omega)|\leq
 \sum_{i:a_i\ne0}|\log|a_i|_\omega|.
\]
Every summand is integrable by the adelic-structure axiom; their finite
sum dominates the absolute value of the measurable local height.
This proves integrability. The zero tuple gives the zero function.
\end{proof}

\begin{definition}[Height of a tuple]\label{def:tuple-height}
\source{chen-moriwaki-adelic:page-0215}
\uses{lem:local-height-integrable}
\lean{ArakelovGeometry.AdelicStructure.logHeight,ArakelovGeometry.AdelicStructure.logHeight_zero}
\leanok
Define $h_S(a)=\int_\Omega\ell_S(a,\omega)\,d\nu(\omega)$.
\end{definition}

\begin{lemma}[Effect of scaling]\label{lem:height-scaling}
\uses{def:tuple-height}
\lean{ArakelovGeometry.AdelicStructure.localHeight_smul,ArakelovGeometry.AdelicStructure.log_localHeight_smul,ArakelovGeometry.AdelicStructure.logHeight_smul}
\leanok
If $a\ne0$ and $t\in K^\times$, then
\[
 \ell_S(ta,\omega)=\log|t|_\omega+\ell_S(a,\omega),\qquad
 h_S(ta)=\int_\Omega\log|t|_\omega\,d\nu(\omega)+h_S(a).
\]
\end{lemma}
\begin{proof}
Multiplicativity and $|t|_\omega>0$ give
$M_S(ta,\omega)=|t|_\omega M_S(a,\omega)$.
Both factors are positive; take their logarithms. The logarithmic
evaluation and local height are integrable, so integration is additive.
\end{proof}

\begin{proposition}[Scalar invariance]\label{prop:height-scalar-invariant}
\source{chen-moriwaki-adelic:page-0215}
\uses{lem:height-scaling,def:adelic-structure}
\lean{ArakelovGeometry.AdelicStructure.logHeight_smul_eq}
\leanok
If $S$ is proper, $a\ne0$ and $t\ne0$, then $h_S(ta)=h_S(a)$.
\end{proposition}
\begin{proof}
The integral of $\log|t|$ in Lemma~\ref{lem:height-scaling} is zero
by the product formula.
\end{proof}

\begin{lemma}[Nonnegativity]\label{lem:height-nonnegative}
\uses{def:tuple-height,def:adelic-structure}
\lean{ArakelovGeometry.AdelicStructure.logHeight_nonneg}
\leanok
If $S$ is proper, then $h_S(a)\geq0$ for every tuple $a$.
\end{lemma}
\begin{proof}
For $a\ne0$ choose $j$ with $a_j\ne0$. The pointwise inequality
$\log|a_j|_\omega\leq\ell_S(a,\omega)$ integrates to
$0\leq h_S(a)$ by integrability and the product formula.
The height of the zero tuple is zero by definition.
\end{proof}

\begin{definition}[Projective height]\label{def:projective-height}
\source{chen-moriwaki-adelic:page-0216}
\uses{prop:height-scalar-invariant,lem:height-nonnegative}
\lean{ArakelovGeometry.AdelicStructure.projectiveHeight,ArakelovGeometry.AdelicStructure.projectiveHeight_mk,ArakelovGeometry.AdelicStructure.projectiveHeight_nonneg}
\leanok
For proper $S$ define $h_S:\mathbb P(K^I)\to\mathbb R$ by
$h_S([a])=h_S(a)$ for any nonzero representative $a$. Its values are nonnegative.
\end{definition}
\begin{proof}
Two nonzero representatives of a projective point differ by a nonzero
scalar, so scalar invariance proves independence of the representative.
The values are nonnegative by Lemma~\ref{lem:height-nonnegative}.
\end{proof}

\begin{definition}[Affine height]\label{def:affine-height}
\uses{def:tuple-height,def:projective-height}
\lean{ArakelovGeometry.AdelicStructure.localHeight_pair,ArakelovGeometry.AdelicStructure.affineHeight,ArakelovGeometry.AdelicStructure.affineHeight_eq_integral,ArakelovGeometry.AdelicStructure.integrable_log_max_one,ArakelovGeometry.AdelicStructure.affineHeight_eq_projectiveHeight}
\leanok
For $b\in K$, write
$h_S^{\mathrm{aff}}(b)=\int_\Omega\log\max(1,|b|_\omega)\,d\nu(\omega)$.
Then $h_S^{\mathrm{aff}}(b)=h_S(1,b)$ and, if $S$ is proper,
$h_S^{\mathrm{aff}}(b)=h_S([1:b])$.
\end{definition}
\begin{proof}
The coordinate $1$ has absolute value one at every place, so the
two local integrands agree. The projective equality follows from the
definition of projective height.
\end{proof}

\begin{lemma}[Affine heights and powers]\label{lem:affine-height-power}
\uses{def:affine-height}
\lean{ArakelovGeometry.AdelicStructure.affineHeight_nonneg,ArakelovGeometry.AdelicStructure.affineHeight_zero,ArakelovGeometry.AdelicStructure.affineHeight_one,ArakelovGeometry.AdelicStructure.affineHeight_neg,ArakelovGeometry.AdelicStructure.affineHeight_pow}
\leanok
Affine heights are nonnegative, $h_S^{\mathrm{aff}}(0)=0$, and
$h_S^{\mathrm{aff}}(b^m)=m h_S^{\mathrm{aff}}(b)$ for every $m\geq0$.
Also $h_S^{\mathrm{aff}}(-b)=h_S^{\mathrm{aff}}(b)$.
\end{lemma}
\begin{proof}
The local maximum is at least one, so its logarithm is nonnegative.
For $b=0$ the maximum is one. For $m=0$ both sides of the power
identity are zero. For $m>0$, multiplicativity and monotonicity of
$x\mapsto x^m$ on $[0,\infty)$ give
$\max(1,|b^m|)=\max(1,|b|)^m$; its logarithm is
$m\log\max(1,|b|)$. Integrating proves the formula.
Finally $|-b|=|b|$ at every place.
\end{proof}

\begin{proposition}[Bounds on affine chart coordinates]
\label{prop:height-coordinate-bound}
\uses{prop:height-scalar-invariant,def:affine-height}
\lean{ArakelovGeometry.AdelicStructure.affineHeight_le_logHeight_of_eq_one,ArakelovGeometry.AdelicStructure.affineHeight_div_le_logHeight}
\leanok
If $S$ is proper, $a_j\ne0$, and $i\in I$, then
$h_S^{\mathrm{aff}}(a_i/a_j)\leq h_S(a)$.
\end{proposition}
\begin{proof}
Set $b=a_j^{-1}a$, so $b_j=1$ and $b_i=a_i/a_j$.
Consequently $\max(1,|b_i|_\omega)\leq M_S(b,\omega)$.
The two positive expressions have integrable logarithms. Integrate
their logarithmic inequality and use $h_S(b)=h_S(a)$.
\end{proof}

\begin{proposition}[Invariance under algebraic extension]
\label{prop:height-algebraic-extension}
\uses{def:tuple-height,def:affine-height,def:algebraic-extension,prop:algebraic-base-pullback}
\lean{ArakelovGeometry.AdelicStructure.localHeight_algebraicExtension,ArakelovGeometry.AdelicStructure.logHeight_algebraicExtension,ArakelovGeometry.AdelicStructure.affineHeight_algebraicExtension}
\leanok
If $L/K$ is algebraic and $a\in K^I$, then
$h_{S\otimes_K L}(a)=h_S(a)$. In particular affine heights of
elements of $K$ are unchanged by algebraic extension.
\end{proposition}
\begin{proof}
For $x\in\Omega_L$ and $b\in K$, one has
$|b|_x=|b|_{\pi_{L/K}(x)}$. Finite maxima and their logarithms
therefore pull back along $\pi_{L/K}$. The base local height is
integrable, so Proposition~\ref{prop:algebraic-base-pullback} gives
equality of its integrals. Apply this to $(1,b)$ for the last assertion.
\end{proof}

\begin{definition}[Height over an algebraic closure]
\label{def:absolute-projective-height}
\source{chen-moriwaki-adelic:page-0215}
\dcref{Definition 3.5.1}
\uses{def:projective-height,def:algebraic-extension,prop:algebraic-proper}
\lean{ArakelovGeometry.AdelicStructure.absoluteHeight,ArakelovGeometry.AdelicStructure.absoluteHeight_mk,ArakelovGeometry.AdelicStructure.absoluteHeight_nonneg}
\leanok
Let $S$ be proper and let $\overline K$ be an algebraic closure of $K$.
The height of $x\in\mathbb P^n(\overline K)$ with respect to $S$ is
$h_S(x):=h_{S\otimes_K\overline K}(x)$.
Equivalently, for any nonzero representative $(a_0,\ldots,a_n)$,
\[
 h_S(x)=\int_{\Omega_{\overline K}}
 \log\max_{0\leq i\leq n}|a_i|_\chi\,d\nu_{\overline K}(\chi).
\]
\end{definition}
\begin{proof}
The extension is proper by Proposition~\ref{prop:algebraic-proper}.
Its projective height is well-defined by
Definition~\ref{def:projective-height}, with precisely the displayed integral.
\end{proof}

\begin{definition}[Northcott property]\label{def:northcott-property}
\lean{ArakelovGeometry.AdelicStructure.HasNorthcott,ArakelovGeometry.AdelicCurve.HasNorthcott,ArakelovGeometry.AdelicStructure.hasNorthcott_iff,ArakelovGeometry.AdelicStructure.absoluteHeight_mk_one}
\source{chen-moriwaki-adelic:page-0216}
\dcref{Definition 3.5.2}
\uses{def:affine-height,lem:affine-height-power,prop:height-algebraic-extension,def:absolute-projective-height}
\leanok
A proper adelic curve $S$ has the Northcott property if
$\{b\in K:h_S^{\mathrm{aff}}(b)\leq C\}$ is finite for every real
$C\geq0$. Equivalently, use the heights $h_S([1:b])$ over $\overline K$.
The same finiteness then holds for every real $C$.
\end{definition}
\begin{proof}
The two height functions agree by algebraic-extension invariance.
For $C<0$, nonnegativity makes the set empty, which is finite.
\end{proof}
